\documentclass[11pt]{article}

\usepackage{amsmath,amssymb,amsthm,mathtools}
\usepackage[margin=1.1in]{geometry}
\usepackage{hyperref}

\newtheorem{theorem}{Theorem}
\newtheorem{lemma}{Lemma}
\newtheorem{corollary}{Corollary}
\newtheorem{remark}{Remark}
\newtheorem{definition}{Definition}

\DeclareMathOperator{\Hom}{Hom}
\newcommand{\K}{K}
\newcommand{\bbE}{\mathbb{E}}
\newcommand{\cC}{\mathcal{C}}
\newcommand{\cH}{\mathcal{H}}
\newcommand{\cT}{\mathcal{T}}

\title{A Chromatic-Polynomial Lower Bound for Pure Chordal Graphs}
\author{}
\date{}

\begin{document}

\maketitle

\section{Motivation}

Let \(H\) be a finite graph and let \(W\colon [0,1]^2\to[0,1]\) be a graphon.
We write
\[
t(H,W)
=
\int_{[0,1]^{V(H)}}
\prod_{ij\in E(H)} W(x_i,x_j)\,
\prod_{i\in V(H)} dx_i
\]
for the homomorphism density of \(H\) in \(W\). Let
\[
p=t(\K_2,W)
\]
be the edge density of \(W\).

For the balanced complete \(k\)-partite graphon \(W_k\), a random map
\[
\phi\colon V(H)\to [k]
\]
is a homomorphism into \(W_k\) precisely when it is a proper \(k\)-colouring of
\(H\). Hence
\[
t(H,W_k)=\frac{\chi_H(k)}{k^{v(H)}}.
\]
Since
\[
t(\K_2,W_k)=1-\frac1k,
\]
if \(p_k=1-1/k\), then \(k=1/(1-p_k)\), and therefore
\[
t(H,W_k)
=
(1-p_k)^{v(H)}
\chi_H\!\left(\frac1{1-p_k}\right).
\]
This motivates the polynomial
\[
\Phi_H(p)
:=
(1-p)^{v(H)}
\chi_H\!\left(\frac1{1-p}\right).
\]

\begin{remark}[Normalisation]
The factor \((1-p)^{v(H)}\), rather than \(p^{-v(H)}\), is forced by the identity
\(t(H,W_k)=\chi_H(k)/k^{v(H)}\). Thus the natural conjectural inequality is
\[
t(H,W)\ge \Phi_H(p).
\]
\end{remark}

The clique density theorem of Reiher proves the analogous optimal lower bound
for cliques. This suggests that balanced multipartite graphons should also be
relevant for other non-bipartite graphs in suitable Turán-type edge-density
ranges. The purpose of this note is to prove such an inequality for a large class:
pure chordal graphs.

\section{Pure chordal graphs and the polynomial \(\Phi_H\)}

\begin{definition}
A graph \(H\) is \emph{chordal} if every cycle of length at least \(4\) has a
chord. Equivalently, its maximal cliques admit a \emph{clique tree}: a tree \(T\)
whose vertices are the maximal cliques \(C_1,\dots,C_m\) of \(H\), such that for
every vertex \(x\in V(H)\), the set of cliques \(C_i\) containing \(x\) induces a
connected subtree of \(T\).
\end{definition}

\begin{definition}
A connected chordal graph \(H\) is called \emph{\(r\)-pure chordal} if every
maximal clique of \(H\) has size \(r\). We assume \(r\ge 3\).
\end{definition}

Let \(H\) be connected \(r\)-pure chordal. Let
\[
C_1,\dots,C_m
\]
be its maximal cliques, and let \(T\) be a clique tree on these cliques. For
\(e=ij\in E(T)\), define the separator
\[
S_e=C_i\cap C_j,
\qquad
s_e=|S_e|.
\]
For \(s\ge 0\), define the polynomial
\[
A_s(p)
:=
\prod_{a=0}^{s-1}\bigl(1-a(1-p)\bigr),
\]
with the convention \(A_0(p)=A_1(p)=1\). Equivalently,
\[
A_s(p)
=
(1-p)^s\left(\frac1{1-p}\right)_s,
\]
where
\[
(x)_s=x(x-1)\cdots(x-s+1).
\]

\begin{lemma}[Chromatic polynomial of a pure chordal graph]
Let \(H\) be connected \(r\)-pure chordal with clique tree \(T\). Then
\[
\chi_H(x)
=
\frac{(x)_r^m}{\prod_{e\in E(T)}(x)_{s_e}}.
\]
Consequently,
\[
\Phi_H(p)
=
(1-p)^{v(H)}
\chi_H\!\left(\frac1{1-p}\right)
=
\frac{A_r(p)^m}{\prod_{e\in E(T)}A_{s_e}(p)}.
\]
\end{lemma}

\begin{proof}
Root the clique tree \(T\). The root clique can be coloured in \((x)_r\) ways.
Whenever a new maximal clique \(C_j\) is attached to its parent clique \(C_i\),
the colours on the separator
\[
S_e=C_i\cap C_j
\]
have already been chosen and are pairwise distinct. The remaining
\(r-s_e\) vertices of \(C_j\) may therefore be coloured in
\[
(x-s_e)_{r-s_e}
=
\frac{(x)_r}{(x)_{s_e}}
\]
ways. Multiplying over the clique tree gives
\[
\chi_H(x)
=
(x)_r
\prod_{e\in E(T)}
\frac{(x)_r}{(x)_{s_e}}
=
\frac{(x)_r^m}{\prod_{e\in E(T)}(x)_{s_e}}.
\]
Also, the clique-tree inclusion-exclusion identity gives
\[
v(H)=mr-\sum_{e\in E(T)}s_e.
\]
Substituting \(x=1/(1-p)\) and multiplying by \((1-p)^{v(H)}\) yields
\[
\Phi_H(p)
=
\frac{A_r(p)^m}{\prod_{e\in E(T)}A_{s_e}(p)}.
\]
\end{proof}

\section{The main inequality}

\begin{theorem}[Chromatic-polynomial lower bound for pure chordal graphs]
Let \(H\) be a connected \(r\)-pure chordal graph with \(r\ge 3\). Let \(W\) be
a graphon with edge density
\[
p=t(\K_2,W).
\]
If
\[
p\ge 1-\frac1{r-1},
\]
then
\[
\boxed{
t(H,W)
\ge
\Phi_H(p)
=
(1-p)^{v(H)}
\chi_H\!\left(\frac1{1-p}\right).
}
\]
Equivalently, if \(C_1,\dots,C_m\) are the maximal \(r\)-cliques of \(H\) and
\(s_e=|C_i\cap C_j|\) are the separator sizes in a clique tree, then
\[
t(H,W)
\ge
\frac{A_r(p)^m}{\prod_{e\in E(T)}A_{s_e}(p)}.
\]
\end{theorem}

\begin{proof}
Write
\[
t_s:=t(\K_s,W).
\]

The proof has two ingredients.

First, we use the Moon--Moser clique-ratio inequality. In graphon form, it says
that for \(2\le j\le r\),
\[
\frac{t_j}{t_{j-1}}
\ge
1-(j-1)(1-p),
\]
whenever the right-hand side is non-negative. Under the assumption
\[
p\ge 1-\frac1{r-1},
\]
this is non-negative for every \(2\le j\le r\). Therefore, for \(0\le s\le r\),
\[
\frac{t_r}{t_s}
=
\prod_{j=s+1}^{r}\frac{t_j}{t_{j-1}}
\ge
\prod_{j=s+1}^{r}\bigl(1-(j-1)(1-p)\bigr)
=
\frac{A_r(p)}{A_s(p)}.
\]
In particular,
\[
t_r\ge A_r(p).
\]

Second, we prove a clique-tree gluing inequality. For every graphon \(W\),
\[
t(H,W)
\ge
\frac{t_r^m}{\prod_{e\in E(T)}t_{s_e}}.
\]
It is enough to prove this first for finite graphs and then pass to graphons by
the usual approximation argument.

Let \(G\) be a finite graph on \(n\) vertices. Put
\[
h_s:=\hom(\K_s,G).
\]
We claim that
\[
\hom(H,G)
\ge
\frac{h_r^m}{\prod_{e\in E(T)}h_{s_e}}.
\]
Root the clique tree \(T\). We construct a random homomorphism from \(H\) to
\(G\). First choose the image of the root \(r\)-clique uniformly from all
homomorphisms \(\K_r\to G\). Then proceed away from the root. Suppose a clique
\(C_j\) is attached to its parent \(C_i\) along the separator
\[
S_e=C_i\cap C_j.
\]
Once the image of \(S_e\) is known, choose the remaining vertices of \(C_j\)
according to the conditional distribution induced by a uniformly random
homomorphism \(\K_r\to G\), conditioned on its restriction to \(S_e\).

The running-intersection property of the clique tree guarantees that the choices
are consistent and produce a homomorphism \(H\to G\). Let \(X\) be the resulting
random homomorphism. Since \(X\) is supported on \(\Hom(H,G)\),
\[
H(X)\le \log \hom(H,G),
\]
where \(H(\cdot)\) denotes Shannon entropy.

For the root clique, the entropy contribution is \(\log h_r\). For every
oriented clique-tree edge \(e\), the conditional entropy of the newly added
vertices is
\[
\log h_r - H(Y_{S_e}),
\]
where \(Y\) is a uniformly random homomorphism \(\K_r\to G\), restricted to
\(S_e\). Since \(Y_{S_e}\) is supported on \(\Hom(\K_{s_e},G)\),
\[
H(Y_{S_e})\le \log h_{s_e}.
\]
Hence
\[
\log \hom(H,G)
\ge
H(X)
\ge
m\log h_r-\sum_{e\in E(T)}\log h_{s_e}.
\]
Exponentiating gives
\[
\hom(H,G)
\ge
\frac{h_r^m}{\prod_{e\in E(T)}h_{s_e}}.
\]
Since
\[
h_s=n^s t(\K_s,G)
\]
and
\[
v(H)=mr-\sum_{e\in E(T)}s_e,
\]
normalising by \(n^{v(H)}\) gives
\[
t(H,G)
\ge
\frac{t_r^m}{\prod_{e\in E(T)}t_{s_e}}.
\]
Passing to graphon limits yields the same inequality for \(W\).

Combining the two ingredients,
\[
t(H,W)
\ge
\frac{t_r^m}{\prod_{e\in E(T)}t_{s_e}}
=
t_r\prod_{e\in E(T)}\frac{t_r}{t_{s_e}}.
\]
Using the Moon--Moser ratio bound,
\[
t_r\prod_{e\in E(T)}\frac{t_r}{t_{s_e}}
\ge
A_r(p)\prod_{e\in E(T)}\frac{A_r(p)}{A_{s_e}(p)}
=
\frac{A_r(p)^m}{\prod_{e\in E(T)}A_{s_e}(p)}.
\]
By the chromatic-polynomial factorisation above, this is exactly
\[
(1-p)^{v(H)}
\chi_H\!\left(\frac1{1-p}\right).
\]
This completes the proof.
\end{proof}

\section{Concrete graph families}

\begin{corollary}[\((r-1)\)-trees]
Let \(H\) be obtained from one \(\K_r\) by repeatedly gluing a new \(\K_r\) along
an existing \(\K_{r-1}\). Suppose \(H\) has \(m\) maximal \(r\)-cliques. Then,
for every graphon \(W\) with
\[
p=t(\K_2,W)\ge 1-\frac1{r-1},
\]
one has
\[
t(H,W)
\ge
A_r(p)\bigl(1-(r-1)(1-p)\bigr)^{m-1}.
\]
\end{corollary}

\begin{proof}
Every separator has size \(r-1\), so
\[
\Phi_H(p)
=
\frac{A_r(p)^m}{A_{r-1}(p)^{m-1}}
=
A_r(p)
\left(\frac{A_r(p)}{A_{r-1}(p)}\right)^{m-1}.
\]
Since
\[
\frac{A_r(p)}{A_{r-1}(p)}
=
1-(r-1)(1-p),
\]
the result follows.
\end{proof}

\begin{corollary}[Generalised books]
Let
\[
B_{r,m}:=\K_{r-1}\vee \overline{\K_m},
\]
namely \(m\) copies of \(\K_r\) sharing a common \(\K_{r-1}\). Then
\[
t(B_{r,m},W)
\ge
A_{r-1}(p)\bigl(1-(r-1)(1-p)\bigr)^m
\]
for every graphon \(W\) with
\[
p\ge 1-\frac1{r-1}.
\]
In particular, for the usual book of \(m\) triangles,
\[
t(B_{3,m},W)\ge p(2p-1)^m,
\qquad p\ge \frac12.
\]
\end{corollary}

\begin{corollary}[Clique windmills]
Let \(F_{r,s,m}\) be the graph formed from \(m\) copies of \(\K_r\), all sharing
the same \(\K_s\), with no other overlap. Then
\[
t(F_{r,s,m},W)
\ge
A_s(p)\left(\frac{A_r(p)}{A_s(p)}\right)^m
\]
whenever
\[
p\ge 1-\frac1{r-1}.
\]
For \(r=3\) and \(s=1\), this gives the friendship graph bound
\[
t(F_{3,1,m},W)
\ge
\bigl[p(2p-1)\bigr]^m,
\qquad p\ge \frac12.
\]
\end{corollary}

\begin{corollary}[Pure triangle-trees]
Let \(H\) be a connected chordal graph whose maximal cliques are all triangles.
Suppose \(H\) has \(m\) triangles and that, in a clique tree, exactly \(a\) of the
separators are edges. Then
\[
t(H,W)
\ge
p^{m-a}(2p-1)^m,
\qquad p\ge \frac12.
\]
\end{corollary}

\begin{proof}
Here \(r=3\),
\[
A_3(p)=p(2p-1),
\qquad
A_2(p)=p,
\qquad
A_1(p)=1.
\]
Thus
\[
\Phi_H(p)
=
\frac{A_3(p)^m}{A_2(p)^a}
=
p^{m-a}(2p-1)^m.
\]
\end{proof}

\section{References}

The proof uses the following external ingredients.

\begin{thebibliography}{9}

\bibitem{Reiher}
C. Reiher,
\emph{The clique density theorem},
Annals of Mathematics \textbf{184} (2016), 683--707.

\bibitem{MoonMoser}
J. W. Moon and L. Moser,
\emph{On a problem of Turán},
Magyar Tud. Akad. Mat. Kutató Int. Közl. \textbf{7} (1962), 283--286.

\bibitem{Agnarsson}
G. Agnarsson,
\emph{On chordal graphs and their chromatic polynomials},
Mathematica Scandinavica \textbf{93} (2003), 240--246.

\bibitem{Zhao}
Y. Zhao,
\emph{Graph Theory and Additive Combinatorics: Exploring Structure and Randomness},
Cambridge University Press, 2023.
See especially the chapter on graph homomorphism inequalities and entropy methods.

\end{thebibliography}

\end{document}